% =====================================================================
% Introduction for EMMET-momentum-DP paper.
% Length: ~1.5 pages in two-column Elsevier format.
% Five subsections in prose (not numbered): problem, gap, contribution,
% results, organization.
% =====================================================================

\section{Introduction}
\label{sec:introduction}

Loss-sensitive workloads---real-time video, financial trading,
telesurgery, autonomous-vehicle telemetry---place strict requirements
on the tail behavior of network paths. A single dropped packet
can stall a TCP flow, jitter a video frame, or invalidate a
deterministic-latency contract. As traffic patterns become
increasingly bursty and topologies host an ever-greater density of
flows, the central problem of adaptive routing remains
unsolved at the algorithmic level: \emph{given a transient view of
edge utilization, choose paths that not only avoid current
congestion but also avoid \textbf{contributing} to future congestion}.

Existing adaptive routers respond to this challenge with edge-level
state. Backpressure routing~\cite{tassiulas1992stability,
ji2013delay} uses queue gradients to drive packet
forward-progress decisions. Load-aware shortest-path (LASP) routing weights
shortest-path costs by edge load, biasing routes away from
busy links. Reinforcement-learning routers~\cite{boyan1994packet,
valadarsky2017learning} learn policies that incorporate richer
features but ultimately collapse to per-edge decisions
indistinguishable across packets. All of these approaches share a
structural property: at any instant, two packets identical in
source, destination and arrival time observe the same network
state, and therefore receive the same routing decision. There is no
mechanism by which one packet's traversal informs the routing of a
sibling packet that has already been issued. Per-packet
\emph{history}---the trajectory each packet has experienced and
the work it has done against the network---is invisible to the
router.

We argue that this missing dimension matters. When a network is
congested, two packets headed for the same destination should not
both follow the same hot path; the router should distribute them
across alternatives even when the global state alone offers no
basis for the distinction. Achieving this without a centralized
controller, without out-of-band coordination, and without violating
delivery guarantees is the question we address.

\paragraph{Contribution.}
We propose \emph{mass-aware routing}, a load-aware routing scheme in
which each packet carries an internal scalar state---its
\emph{mass}---that evolves with the packet's trajectory. The mass
starts at unity and grows multiplicatively whenever the packet
traverses a congested edge:
\begin{equation}
m_{\text{out}} \;=\; \min\!\Big(m_{\text{in}}\cdot
                       \big(1 + \kappa\,\rho(u,v)\big),\; m_{\max}\Big),
\label{eq:mass-update-intro}
\end{equation}
where $\rho(u,v)$ is the local congestion of edge $(u,v)$ and
$\kappa$ is a growth coefficient. At every routing decision, the
mass scales the cost of candidate edges:
\(\delta(u,v) \;=\; m \cdot \Phi(u,v)\), where $\Phi$ is an integrated
edge potential combining latency, congestion, and a thermal-memory
term tracking past loss events. A packet that has accumulated
mass through a congested neighborhood is, by construction, more
expensive to route through any further congested edges, and
therefore the router---given the same destination---steers it
toward less-loaded alternatives.

We instantiate this principle as a constrained dynamic-programming
router (\emph{EMMET-DP}) over the joint state space
$\{\text{node}\}\times\{\text{hops}\}\times\{\text{mass-bucket}\}$,
with hop budget $H = \lceil\alpha_b\cdot h_{\text{sp}}\rceil$
relative to the unweighted shortest path. The DP guarantees that
a feasible path within budget is always found when one exists,
eliminating the dead-end failure modes of greedy adaptive routers.
The complete algorithm is given in
Algorithm~\ref{alg:massaware-dp}.

\paragraph{Why per-packet state matters.}
Two structural properties distinguish mass-aware routing from
prior load-aware schemes. First, two packets identical in
$(s,t,\text{arrival time})$ but differing in trajectory thus far
receive \emph{different} routing recommendations: the heavier
packet is steered toward cooler regions of the graph, the lighter
packet may take the more direct route. This produces emergent load
balancing without any central coordinator. Second, because the
mass is a property of the packet itself rather than the network,
the mechanism degrades gracefully under workloads that conventional
load-aware routers handle poorly: when a sequence of packets shares
the same source and destination, mass-aware routing spreads them
across paths in proportion to instantaneous network state, whereas
LASP-class algorithms send them all down the same locally-optimal
edge.

\paragraph{Results.}
On the GÉANT topology with $100$ random traffic seeds, mass-aware
DP reduces congestion losses by $60.2\%$ (95\% CI $[-64.1, -36.7]$\%)
relative to LASP-aug, a Dijkstra baseline using the same integrated
potential function (paired $t = 5.21$, Wilcoxon $p = 8.4\times 10^{-9}$,
effect size $r = 0.73$, Cohen's $d = 0.52$). Across all 26 evaluated
scenarios ($n = 2{,}400$ paired seeds), the pooled paired test is
$t = 6.90$ and Wilcoxon $p = 8.6\times 10^{-16}$. The
improvement holds with negligible additional capacity consumption
per delivered packet ($+0.04$ hops on average). The 26 scenarios
span five topology families: Erd\H{o}s--R\'enyi
($18$ configurations), scale-free Barab\'asi--Albert,
small-world Watts--Strogatz, $7\times 7$ regular lattice, and the
real GÉANT and Abilene backbones. Crucially, the improvement
\emph{vanishes} in uncongested regimes---when the network has no
congestion to disambiguate, mass-aware routing reduces to the
underlying shortest-path baseline and gains nothing. This
selective effectiveness, rather than a uniform improvement, is
the signature of a mechanism that extracts genuinely new
information from per-packet state.

Mass-aware DP is approximately $8.8\times$ slower than LASP-aug
per routing decision, situating it firmly in the controller-plane
or traffic-engineering regime rather than per-packet forwarding.
We discuss the implications of this trade-off, and possible
hardware-accelerated variants, in Section~\ref{sec:discussion}.

\paragraph{Organization.}
Section~\ref{sec:related} surveys load-aware and stateful routing
literature.  Section~\ref{sec:model} formalizes the network model,
the integrated edge potential, and the mass-update dynamics.
Section~\ref{sec:algorithm} presents the EMMET-DP algorithm
and analyzes its complexity. Section~\ref{sec:experiments} details
the evaluation methodology, including baselines, topologies and
statistical protocol. Section~\ref{sec:results} reports per-scenario
results and ablations. Section~\ref{sec:discussion} discusses
limitations, deployment regimes, and connections to prior work.

